% Generated draft: requires mathematical and source review.
\documentclass[11pt]{article}
\usepackage[margin=25mm]{geometry}
\usepackage{fontspec}
\setmainfont{Latin Modern Roman}
\usepackage{amsmath,amssymb,mathtools,unicode-math}
\setmathfont{Latin Modern Math}
\usepackage{enumitem,needspace,xurl,hyperref}
\setlength{\parindent}{0pt}
\setlength{\parskip}{5pt}
\setlength{\emergencystretch}{3em}
\newcommand{\N}{\mathbb{N}}
\newcommand{\Z}{\mathbb{Z}}
\newcommand{\Q}{\mathbb{Q}}
\newcommand{\R}{\mathbb{R}}
\newcommand{\C}{\mathbb{C}}
\newcommand{\F}{\mathbb{F}}
\newcommand{\A}{\mathbb{A}}
\newcommand{\PP}{\mathbb{P}}
\newcommand{\E}{\mathbb{E}}
\newcommand{\eps}{\varepsilon}
\newcommand{\dd}{\,\mathrm{d}}
\begin{document}
\section*{UnsolvedMath ID: 20000544}
\textbf{Source problem code:} catalogue entry
\subsection{Definitions and mathematical statement}
The primary AIM question asks when pure states of a MASA $A$ extend uniquely to states of a $C^*$-algebra $B=I\oplus A$, with $I$ an ideal. The report interprets this as an internal linear splitting $B=I\dotplus A$, with $I\cap A=0$, rather than a $C^*$ direct product. In its unital formulation, $A$ contains the unit of $B$, is abelian and satisfies $A'\cap B=A$. A state is a positive linear functional of norm one; a pure state of $A$ is a character $\phi$. The original does not specify unitality or the splitting convention.

The report defines the positive eigenfunctional cone
\[
 \mathcal C_\phi=\{\tau\in I^*_+:\tau(ai)=\phi(a)\tau(i),\
 \tau(ia)=\tau(i)\phi(a)\ \text{for all }a\in A,i\in I\}.
\]
It claims the states extending $\phi$ are exactly
\[
 \psi_\tau(i+a)=\tau(i)+\phi(a),\qquad
 \tau\in\mathcal C_\phi,\quad\|\tau\|\leq1.
\]
Thus uniqueness is equivalent to $\mathcal C_\phi=\{0\}$, the unique extension then factoring through $B/I\cong A$. If $I\cong\mathcal K(H)$ and $A$ is a MASA, the report claims this cone vanishes for every character.

In contrast it reports a separable unital example with
\[
 I\cong C_0((0,1],\mathcal K(\C\oplus L^2[0,1])),\qquad A\cong C([0,1])
\]
a MASA in an internal split $B$, such that evaluation at $1/2$ has two distinct pure extensions. Its stated mechanism is an eigenvector appearing only in the endpoint fiber. The embedding of $A$, field of representations and formula for $B$ are not supplied, so the claimed explicit example cannot be reconstructed without further data. It shows the report's asserted insufficiency of MASA plus linear splitting, without contradicting the compact-operator-ideal case.

The complete record also appends a participant agenda, retained here as attributed context rather than additional hypotheses or resolved results:
\begin{itemize}
\item Cynthia Vinzant asks for constructive Weaver/vector-partitioning and paving methods, including a possible two-player communication-complexity obstruction. Greg Knese asks for intrinsic Feichtinger proofs, multivariable interlacing and quantitative root control by stability-preserving operators. Petter Branden asks for constructive matching lower bounds, improved examples, continuous matroid consequences and hyperbolic frame/polynomial extensions. The notes mention a necessary $\sqrt\varepsilon$ scale and real paving figures $1/\varepsilon^2$ and $6/\varepsilon^2$, without complete statements, and refer to an unidentified ``Pete C.'' example and an unspecified new preprint.
\item Marcin Bownik asks for the minimum and construction of Riesz-sequence partitions of Parseval frames whose vector norms are bounded below, especially exponential frames on measurable subsets of the torus. Darrin Speegle asks whether partition counts depend on measure or finer geometry; specifically, whether an $E\subset[0,1]$ of measure $1/2$ can have its restricted integer-exponential frame fail every partition into two Riesz sequences. The displayed exponential formula is malformed; the usual restricted functions $e^{2\pi int}\chi_E(t)$ are only an explicit interpretation. The agenda asks for constructive arithmetic-progression/``synthetic set'' consequences attributed to Bourgain--Tzafriri; that terminology is not defined.
\item Bill Johnson asks the precise operator-norm paving question: for each $1\leq p<\infty$ and $\varepsilon>0$, is there $k=k(\varepsilon,p)$ such that every $n\times n$ zero-diagonal $D$ with $\|D\|_{\ell_p^n\to\ell_p^n}=1$ has a partition $\sigma_1,\ldots,\sigma_k$ of $\{1,\ldots,n\}$ with
\[
 \left\|\sum_jP_{\sigma_j}DP_{\sigma_j}\right\|
 =\max_j\|P_{\sigma_j}DP_{\sigma_j}\|\leq\varepsilon?
\]
These are not Schatten norms. The notes tentatively attribute known $p=1$ results to Bourgain (``IIRC'') and compare them with nonnegative-entry $p=2$ paving; they cite Berman--Halpern--Kaftal--Weiss. They also recall Bourgain--Tzafriri restricted invertibility and ask for dimension-free commutator bounds $D=[B,C]$ with $B$ diagonal of norm one. The required normalization of $D$, the Johnson--Ozawa--Schechtmann PNAS reference and the recorded comparison of $n^\alpha$ with ``poly$(n)$'' are incomplete.
\item Deanna Needell asks for efficient paving algorithms for stochastic optimization and sampling/RIP constructions; Anna Gilbert asks for sharper DFT-RIP/Fourier sampling bounds and use of Rudelson selection. Olaf Mordhorst requests concrete optimality examples for restricted invertibility. Pierre Youssef asks for unweighted sparse balanced identity approximations from $\sum_jc_jx_jx_j^T=I$, $\|x_j\|=1$, $\sum_jc_jx_j=0$: choose $S$ of size comparable to $d$ with $C\sum_{j\in S}x_jx_j^T\approx I$ and $\|\sum_{j\in S}x_j\|\leq B$. The notes mention Srivastava's $4+\varepsilon$ approximation, sections of dimension $(1-\varepsilon)n$, a reference-body distance of order $\sqrt{\dim E}$, and an $n^{5/6}$ Banach--Mazur-to-cube benchmark, with constants and the body $B_1$ unspecified; they suggest revisiting a Batson--S--S approach.
\item Leonid Gurvits discusses a proposed five-spanning-set/Rado--Horn procedure using repeated vectors and identity padding, but leaves its matrix test as ``xxxx''. His zonotope question uses segments $[0,z_i]$, $k>n$, a weighted Minkowski sum and a proposed balanced two-part approximation; the volume-one phrase and $(1/2\pm\varepsilon)$ inclusion have missing normalizations and directions of comparison. Mirko Visontai asks to count good partitions and generate interlacing families. Mihai Putinar asks for a Nevanlinna/phase-shift interpretation. Dan Edidin asks for dimension-free diagonal criteria for $r$-pavability of projections and geometric parameter spaces/counts of good partitions. Pablo Parrilo and Rachel Ward emphasize algorithms. Bernhard Bodmann asks for multivariable perturbation proofs without the polynomial machinery. Adam Marcus asks for portable tools, worst-root models from repeated identity copies with zero padding, weaker largest-root-only interlacing, and rank-two rather than rank-one perturbations. These agendas contain incomplete formulations and no supplied answers.
\end{itemize}
The research report concerns only the primary pure-state question; it does not resolve this appended workshop program.
\subsection{Short English statement}
Characterize unique pure-state extension across an internally split ideal and MASA. The report gives an exact positive-functional criterion, a compact-operator-ideal positive case and an incompletely specified continuous-field counterexample. The appended workshop agenda is preserved as attributed context, including its independent paving question and incomplete quantitative claims.
\subsection{Sources}
\begin{itemize}
\item Local registry: \url{data/unsolvedmath/problems.json}, permanent ID $20000544$, source code \texttt{catalogue entry}; complete primary question, appended participant agenda and research summary. The full prompt, including repeated source fields, was read in chunks after display truncation.
\item AIM workshop ``Beyond Kadison-Singer: paving and consequences'', Question 3, Google document: \url{https://docs.google.com/document/d/1XRw0WyG5PshqboCp99tovCRD1GCK0jXZiNLH_nw618c/edit}. Canonical location \texttt{aim-analysis-notes.json}, \texttt{notes[75]}; original remarks and literature field empty. No document was visited here.
\item Kenneth Berman, Herbert Halpern, Victor Kaftal and Gary Weiss, ``Matrix norm inequalities and the relative Dixmier property'', \textit{Integral Equations and Operator Theory} 11 (1988), no. 1, pp. 28--48, as supplied in the appended notes.
\item Other supplied references are informal: Bourgain--Tzafriri, a Johnson--Ozawa--Schechtmann PNAS paper, MSS/BSS methods, Srivastava and an unidentified hyperbolic-polynomial preprint. Their complete bibliographic data and theorem hypotheses are absent; contributor attributions are retained in the text.
\item Supplied machine-generated research attempt 1: partial result; problem status unknown at the run, registry partially solved; rigor metadata ``proved'', novelty confidence low. The non-peer-reviewed eigenfunctional and counterexample claims require independent verification before citation.
\end{itemize}
\end{document}
